\documentclass{article}
\usepackage{natbib}
\title{Experiments and Entry: A Short Test Manuscript}
\begin{document}
\maketitle

Entrepreneurs increasingly treat their business ideas as hypotheses to be tested \citep{kerr2014}, and investors have begun to finance ventures in stages so that each stage functions as an experiment \citep{nanda2015}. A firm earns profits from a resource only when it has more accurate expectations about the resource's value than other market participants, or when it is lucky \citep{barney1986}.

Experiments can also change what a firm captures. \citet{schmalensee1981} shows that a firm's experiments reveal market opportunities to potential entrants and thereby induce entry. Firms sometimes counter this by preannouncing new products, which deters entry \citep{lieberman1988}. Evidence from a randomized trial shows that a scientific approach to entrepreneurial decision making improves performance mainly by increasing the rate at which entrepreneurs terminate their ideas \citep{camuffo2020}. Taken together, these results suggests that the value of an experiment depend on who else can act on it.

\bibliographystyle{apalike}
\bibliography{refs}
\end{document}
